% ==================== 5 结论 ====================
\section{结\hspace{0.5em}论}

本文以单晶硅为对象，完成了晶体学计算模板的建立、PDF 卡片的构建、晶面消光系数的计算和五个带轴电子衍射的模拟，主要结果如下。

（1）在 Origin 中建立了一套可复用的晶体学计算模板。模板以 A--F 列输入 $a,b,c,\alpha,\beta,\gamma$，以 G 列计算晶胞体积，以 H--M 列计算倒易点阵参数，以 N--S 列输入两组晶面指数，以 T--V 列计算倒易矢量长度 $r_1,r_2$ 和夹角 $\theta$。对 $a=5.43$~\AA 的单晶硅，模板给出 $V=160.103$~\AA$^3$、$a^*=b^*=c^*=0.18416$~\AA$^{-1}$、$\alpha^*=\beta^*=\gamma^*=90^\circ$，与立方晶系的解析结果一致。

（2）在模板基础上增加了 $2\theta$、多重性、$|F|^2$ 和相对强度四列，构建了单晶硅 PDF 卡片。Cu K$\alpha_1$ 条件下最强的衍射线为 111（$2\theta=28.45^\circ$），其次是 220（$47.31^\circ$）、311（$56.13^\circ$）、422（$88.05^\circ$）等；高角反射的强度对原子散射因子 $f(s)$ 和温度因子很敏感，说明 PDF 卡片的强度计算不能只用常数 $f$。

（3）先由 VESTA 导出的 8 个原子坐标和结构因子计算晶面的消光系数，判断允许/消光；再用晶带方程选取基矢，用式 \eqref{eq:xy} 和式 \eqref{eq:hkl-mn} 建立二维点阵，完成了 [001]、[111]、[332]、[334]、[345] 五个带轴的电子衍射模拟。[001] 为正方形点阵（$r_1=r_2=0.18416$~\AA$^{-1}$，$\theta=90^\circ$），[111] 为三角/六方点阵（$r_1=r_2=0.26044$~\AA$^{-1}$，$\theta=120^\circ$），[332]、[334]、[345] 为不同夹角和间距的斜方点阵。五个带轴的基矢叉积分别平行于对应的带轴，验证了几何关系的正确性。

（4）以 110 反射为例验证了金刚石结构的消光规律。110 的 $h,k,l$ 奇偶混合，8 个原子的实部正负抵消，计算强度约为 $4.7\times10^{-29}$，在数值上为零。这说明几何点阵只给出可能出现的位置，实际衍射花样还必须满足结构因子不为零的条件。

通过本文的工作，进一步理解了晶胞参数、倒易点阵、晶带方程、结构因子和 PDF 卡片之间的联系，也熟悉了 Origin 的公式编辑、列引用、角度单位设置和三维绘图操作。模板目前采用运动学近似和常数 $f=14$，后续可以加入 $f(s)$、Debye--Waller 因子、多重散射和 HOLZ 线，使其更接近实验电子衍射花样。
